% SPDX-License-Identifier: MIT
% Copyright (c) 2026 Evgeny Baulin
% Optimization Methods, Seminar 05: theory handout (Russian).
% Generated by checks/05. Second-Order and Quasi-Newton Methods.ipynb from web/05/js/content-*.js and content/students.json. Do not edit by hand.
% Compile with: latexmk -xelatex -interaction=nonstopmode -halt-on-error Theory_ru.tex
\documentclass[11pt,a4paper]{article}

\usepackage{fontspec}
\setmainfont{cmunrm.otf}[BoldFont=cmunbx.otf, ItalicFont=cmunti.otf, BoldItalicFont=cmunbi.otf]
\setsansfont{cmunss.otf}[BoldFont=cmunsx.otf, ItalicFont=cmunsi.otf, BoldItalicFont=cmunso.otf]
\setmonofont{cmuntt.otf}[BoldFont=cmuntb.otf, ItalicFont=cmunit.otf, BoldItalicFont=cmuntx.otf]
\usepackage{polyglossia}
\setdefaultlanguage{russian}
\setotherlanguage{english}
\newfontfamily\cyrillicfont{cmunrm.otf}[Ligatures=TeX, BoldFont=cmunbx.otf, ItalicFont=cmunti.otf, BoldItalicFont=cmunbi.otf]
\newfontfamily\cyrillicfontsf{cmunss.otf}[Ligatures=TeX, BoldFont=cmunsx.otf, ItalicFont=cmunsi.otf, BoldItalicFont=cmunso.otf]
\newfontfamily\cyrillicfonttt{cmuntt.otf}[BoldFont=cmuntb.otf, ItalicFont=cmunit.otf, BoldItalicFont=cmuntx.otf]

\usepackage[a4paper,left=2.3cm,right=2.3cm,top=2.4cm,bottom=2.4cm,headheight=14pt]{geometry}
\usepackage{amsmath,amssymb}
\usepackage{graphicx}
\usepackage{array,booktabs,longtable}
\usepackage{enumitem}
\usepackage{xcolor}
\usepackage[most]{tcolorbox}
\usepackage{titlesec}
\usepackage{fancyhdr}
\usepackage{needspace}
\usepackage[expansion=false]{microtype}
\usepackage[colorlinks=true,linkcolor=hseblue,urlcolor=hseblue,citecolor=hseblue]{hyperref}

\definecolor{hseblue}{RGB}{16,45,105}
\definecolor{hseaccent}{RGB}{198,40,40}
\definecolor{hsegray}{RGB}{108,114,126}
\definecolor{hselight}{RGB}{240,243,248}
\definecolor{hsewarm}{RGB}{253,244,232}

\hypersetup{pdftitle={Методы оптимизации. Семинар 5. Конспект теории}, pdfauthor={Евгений Баулин}}

\titleformat{\section}{\color{hseblue}\Large\bfseries}{\thesection}{0.6em}{}
\titleformat{\subsection}{\color{hseblue}\large\bfseries}{\thesubsection}{0.6em}{}
\titlespacing*{\section}{0pt}{16pt plus 4pt minus 2pt}{8pt plus 2pt}
\titlespacing*{\subsection}{0pt}{12pt plus 3pt minus 2pt}{6pt plus 2pt}
% a heading keeps a few lines under it on its page: a breakable box right after it would otherwise leave it alone
\pretocmd{\section}{\Needspace{7\baselineskip}}{}{\errmessage{cannot patch section}}
\pretocmd{\subsection}{\Needspace{5\baselineskip}}{}{\errmessage{cannot patch subsection}}

\tcbset{handoutbox/.style={enhanced, breakable, sharp corners, boxrule=0pt, left=7pt, right=7pt, top=5pt, bottom=5pt, before skip=8pt plus 2pt, after skip=8pt plus 2pt}}
\newtcolorbox{examplebox}{handoutbox, colback=hselight, colframe=hselight, before upper={\textcolor{hseblue}{\bfseries Пример.}\ }}
\newtcolorbox{whybox}{handoutbox, colback=white, colframe=hseblue, leftrule=2.5pt, top=3pt, bottom=3pt, before upper={\textcolor{hseaccent}{\bfseries Почему.}\ }}
\newtcolorbox{answerbox}{handoutbox, colback=hsewarm, colframe=hsewarm, before upper={\textcolor{hseaccent}{\bfseries Ответ.}\ }}
\newtcolorbox{conventionbox}{handoutbox, colback=hselight, colframe=hseblue, leftrule=2.5pt, before upper={\textcolor{hseblue}{\bfseries Обозначения.}\par\smallskip}}

\newcommand{\handoutfigure}[3]{\par\medskip\noindent\begin{minipage}{\linewidth}\centering
  \includegraphics[width=#1\linewidth]{figures/#2.pdf}\par\smallskip
  \begin{minipage}{0.94\linewidth}\small\textcolor{hseblue}{\bfseries Рисунок.}\ #3\end{minipage}\end{minipage}\par\medskip}
\newcommand{\problemhead}[3]{\par\bigskip\noindent{\color{hseblue}\large\bfseries #1.\ #2}\par\nopagebreak
  \noindent{\small\itshape\color{hsegray}#3}\par\nopagebreak\smallskip}
\newtcolorbox{definitionbox}[1]{handoutbox, colback=white, colframe=hseblue, leftrule=2.5pt, before upper={\textcolor{hseblue}{\bfseries #1}\ }}
\newtcolorbox{theorembox}[1]{handoutbox, colback=hselight, colframe=hseblue, leftrule=2.5pt, before upper={\textcolor{hseblue}{\bfseries #1}\ }}
\newtcolorbox{proofbox}{handoutbox, colback=white, colframe=hsegray, leftrule=1.2pt, top=3pt, bottom=3pt, before upper={\textcolor{hsegray}{\bfseries\itshape Доказательство.}\ }}
\newtcolorbox{notebox}[1]{handoutbox, colback=hselight, colframe=hselight, before upper={\textcolor{hseblue}{\bfseries #1}\par\smallskip}}
\newtcolorbox{trapbox}{handoutbox, colback=hsewarm, colframe=hseaccent, leftrule=2.5pt, before upper={\textcolor{hseaccent}{\bfseries Ловушка.}\ }}
\newtcolorbox{selfcheckbox}{handoutbox, colback=white, colframe=hseaccent, leftrule=2.5pt, top=3pt, bottom=3pt, before upper={\textcolor{hseaccent}{\bfseries Проверьте себя.}\ }}
% a display wider than the line is set as a centred paragraph that may break after a relation
\newcommand{\handoutdisplay}[1]{\sbox0{$\displaystyle #1$}\ifdim\wd0>\linewidth
  \par{\centering\binoppenalty=10000\relpenalty=0\relax$\displaystyle #1$\par}\else\[#1\]\fi}
\newcommand{\selfcheckanswer}{\par\smallskip\textcolor{hsegray}{\bfseries Ответ.}\ }

\setlist[itemize]{itemsep=3pt,topsep=4pt}
\setlist[enumerate]{itemsep=3pt,topsep=4pt}
\setlength{\parindent}{0pt}
\setlength{\parskip}{5pt plus 1pt}
\setlength{\emergencystretch}{2em}
\setcounter{tocdepth}{2}

\pagestyle{fancy}
\fancyhf{}
\fancyhead[L]{\small\color{hsegray}Методы оптимизации}
\fancyhead[R]{\small\color{hsegray}Семинар 5}
\fancyfoot[L]{\small\color{hsegray}Евгений Баулин}
\fancyfoot[C]{\small\thepage}
\fancyfoot[R]{\small\color{hsegray}Лицензия MIT}
\renewcommand{\headrulewidth}{0pt}
\renewcommand{\footrulewidth}{0pt}
\fancypagestyle{plain}{\pagestyle{fancy}}

\begin{document}

\begin{center}
  {\large\color{hsegray}Методы оптимизации\par}
  \vspace{4pt}
  {\large\bfseries\color{hseaccent}Семинар 5\par}
  \vspace{8pt}
  {\LARGE\bfseries\color{hseblue}Методы второго порядка и квазиньютоновские методы\par}
  \vspace{8pt}
  {\Large Конспект теории\par}
  \vspace{12pt}
  НИУ ВШЭ, факультет компьютерных наук\par
  2026–2027 учебный год\par
  \vspace{8pt}
  Лектор курса: Андрей Игнатов\par
  Евгений Баулин\par
  Telegram: \href{https://t.me/tarakan_tuc}{@tarakan\_tuc}\par
\end{center}

\noindent{\color{hseblue}\rule{\linewidth}{0.6pt}}

\tableofcontents

\clearpage


Градиентный спуск использует только наклон и расплачивается за это числом обусловленности. Конспект следует лекции и компьютерному семинару Андрея Игнатова о методах второго порядка: шаг Ньютона и его локальная теория, защита, без которой он не работает вдали от решения, методы Гаусса–Ньютона и Левенберга–Марквардта, BFGS и L-BFGS, стоимость кривизны.

\begin{conventionbox}
$f:\mathbb R^n\to\mathbb R$ дважды непрерывно дифференцируема, $g=\nabla f(x)$, $H=\nabla^2 f(x)$, $e=x-x^*$. $B$ приближает гессиан, $G$ — обратный к нему; $s_k=x_{k+1}-x_k$, $y_k=g_{k+1}-g_k$. $A\succ0$ означает положительную определённость. Константа Армихо $c\in(0,1)$, обычно малая.
\end{conventionbox}

% =====================================================================
\section{Квадратичная модель и шаг Ньютона}\label{sec:model}

\begin{definitionbox}{Определение.}
\textbf{Квадратичная модель} $f$ в точке $x$ — это
\handoutdisplay{
m_x(p)=f(x)+g^\top p+\tfrac{1}{2}p^\top Hp,
}
 она совпадает с $f(x+p)$ с точностью до $o(\|p\|^2)$. Если $H$ обратима, \textbf{направление Ньютона} $p_N$ решает $Hp_N=-g$, а полный шаг Ньютона равен $x^+=x+p_N$.
\end{definitionbox}

$\nabla_p m_x(p)=g+Hp$, так что $p_N$ — стационарная точка модели. При $H\succ0$ это единственный минимум модели; при знаконеопределённой $H$ — её седло, а при $H\prec0$ — максимум.

\begin{whybox}
Градиентный спуск минимизирует модель, в которой $H$ заменён на $\alpha^{-1}I$: круглую чашу. Ньютон использует чашу настоящей формы. В собственном базисе $H\succ0$ шаг Ньютона делит каждую компоненту градиента на свою кривизну: крутым направлениям достаются короткие шаги, пологим — длинные.
\end{whybox}

\begin{theorembox}{Теорема: Один шаг на квадратичной функции.}
Для $f(x)=\tfrac{1}{2}x^\top Qx-b^\top x+c$ с $Q\succ0$ один полный шаг Ньютона из любой точки $x$ попадает в $x^*=Q^{-1}b$.
\end{theorembox}

\begin{proofbox}
$g=Qx-b$ и $H=Q$. Уравнение Ньютона $Qp=-(Qx-b)$ даёт $p=-x+Q^{-1}b=x^*-x$, поэтому $x+p=x^*$.
\end{proofbox}

\begin{examplebox}
Для $Q=\operatorname{diag}(1;\ 100)$ GD с шагом $\alpha=1/L=0{,}01$ уменьшает ошибку вдоль пологой оси в $0{,}99$ раза за шаг, и чтобы уменьшить её в $1000000$ раз, нужно $1375$ шагов. Шаг Ньютона делит две компоненты градиента на их кривизны и заканчивает за один шаг.
\end{examplebox}

\begin{trapbox}
Решайте $Hp=-g$ через разложение Холецкого или LU, \texttt{numpy.linalg.solve}, но никогда как \texttt{inv(H) @ g}. Явная обратная матрица дороже и менее точна, а неудавшееся разложение Холецкого — самая дешёвая проверка того, что $H\succ0$.
\end{trapbox}

\problemhead{A2}{Один шаг на квадратичной функции}{Самостоятельно}
\textbf{Условие.} Пусть $f(x)=\tfrac{1}{2}x^\top Qx-b^\top x$, $Q=\begin{pmatrix} 4 & 1 \\ 1 & 2 \end{pmatrix}$, $b=(1;\ -1)$. Из $x_0=(2;\ 0)$ найдите шаг Ньютона и $x_1$ и покажите, что $x_1$ не зависит от $x_0$. Найдите число обусловленности $\kappa$ матрицы $Q$.

\textbf{Решение.}
\begingroup\sloppy
\begin{enumerate}[label=\textbf{\arabic*.}, leftmargin=*]
  \item $g_0=Qx_0-b=(7;\ 3)$. Из $Qp=-g_0$ получаем $p_0=(-1{,}5714;\ -0{,}7143)$ и $x_1=x_0+p_0=(0{,}4286;\ -0{,}7143)$.
  \item Для любого $x$ уравнение Ньютона имеет вид $Qp=-(Qx-b)$, поэтому $p=Q^{-1}b-x=x^*-x$ и $x+p=x^*$. Действительно, $Qx_1-b=0$.
  \item Собственные значения $Q$ равны $(1{,}5858;\ 4{,}4142)$, значит, $\kappa=2{,}7836$. Градиентный спуск с лучшим постоянным шагом сжимал бы ошибку в $(\kappa-1)/(\kappa+1)=0{,}4714$ раза за шаг.
\end{enumerate}
\endgroup
\begin{answerbox}
$x_1=(0{,}4286;\ -0{,}7143)=Q^{-1}b$ из любой точки; $\kappa=2{,}7836$.
\end{answerbox}

% =====================================================================
\section{Локальная квадратичная сходимость}\label{sec:local}

\begin{theorembox}{Теорема: Квадратичная сходимость.}
Пусть $\nabla f(x^*)=0$, $\nabla^2 f(x^*)\succ0$ и гессиан $M$-липшицев возле $x^*$. Тогда существуют окрестность $x^*$ и $m>0$, такие что полные шаги Ньютона из неё не выходят и
\handoutdisplay{
\|e^+\|\le{M\over 2m}\|e\|^2.
}

\end{theorembox}

\begin{proofbox}
По формуле Ньютона–Лейбница $g(x)=\int_0^1 H(x^*+te)\,e\,dt$. Так как $e^+=e-H(x)^{-1}g(x)$,
\handoutdisplay{
e^+=H(x)^{-1}\int_0^1\bigl[H(x)-H(x^*+te)\bigr]e\,dt.
}
 Возле $x^*$ гессиан остаётся положительно определённым, $\|H(x)^{-1}\|\le1/m$, и $\|H(x)-H(x^*+te)\|\le M(1-t)\|e\|$. Интеграл от $1-t$ даёт множитель $\tfrac{1}{2}$.
\end{proofbox}

Как только $\tfrac{M}{2m}\|e\|<1$, каждый шаг примерно вдвое увеличивает количество правильных первых знаков $x_k$. Важны обе константы: почти вырожденный гессиан в решении (малое $m$) или быстро меняющийся гессиан (большое $M$) сужает область, где это происходит.

\handoutfigure{0.85}{fig_digits}{Ошибка метода Ньютона на $x-\ln x$ по итерациям в сравнении с градиентным спуском.}

\problemhead{A3}{Знаки удваиваются}{Самостоятельно}
\textbf{Условие.} Примените метод Ньютона к $f(x)=x-\ln x$, $x>0$, из $x_0=0{,}5$. Выведите рекуррентное соотношение для $e_k=1-x_k$ и найдите $x_3$. Сколько шагов нужно, чтобы $e_k\le10^{-12}$? При каких $x_0>0$ метод сходится?

\textbf{Решение.}
\begingroup\sloppy
\begin{enumerate}[label=\textbf{\arabic*.}, leftmargin=*]
  \item $x^+=x-(1-1/x)x^2=2x-x^2$. При $e=1-x$: $e^+=1-2x+x^2=(1-x)^2=e^2$, точно и с $C=1$.
  \item Из $x_0=0{,}5$ итерации равны $(0{,}5;\ 0{,}75;\ 0{,}9375;\ 0{,}9961)$, ошибки — $(0{,}5;\ 0{,}25;\ 0{,}0625;\ 0{,}0039)$; $x_3=0{,}9961$.
  \item $e_k=e_0^{2^k}=2^{-2^k}$. Впервые она становится меньше $10^{-12}$ при $k=6$: каждый шаг удваивает число верных двоичных знаков.
  \item Итерация сходится ровно при $|e_0|<1$, то есть при $0< x_0<2$. Из $x_0=2$ она попадает в $0$, за пределы области определения; из большей точки, например из $x_0=2{,}5$, прыгает в $-1{,}25$.
\end{enumerate}
\endgroup
\begin{answerbox}
$e^+=e^2$, $x_3=0{,}9961$, шагов: $6$; сходимость при $0< x_0<2$.
\end{answerbox}

\begin{trapbox}
Для $f(x)=x^4$ в минимуме $f''(0)=0$, и метод Ньютона даёт $x^+=0{,}6667\,x$: линейная сходимость, $35$ шагов на уменьшение в $1000000$ раз. Теорема выше не зря требует $\nabla^2 f(x^*)\succ0$.
\end{trapbox}

% =====================================================================
\section{Когда полный шаг не работает}\label{sec:failure}

Вдоль $p_N$ имеем $g^\top p_N=-g^\top H^{-1}g$. Это отрицательно при $H\succ0$ и может быть положительно при знаконеопределённой $H$. Даже при $H\succ0$ единичный шаг может выйти из области, где модель описывает $f$.

\begin{itemize}
  \item \textbf{Направление вверх.} Знаконеопределённая $H$ может дать $g^\top p_N>0$; линейный поиск вдоль $p_N$ не поможет.
  \item \textbf{Не та стационарная точка.} Метод Ньютона решает $\nabla f=0$ и притягивается к сёдлам и максимумам так же охотно, как к минимумам.
  \item \textbf{Слишком длинный шаг.} Модель верна в $x$ и неверна там, где лежит её минимум.
  \item **Вырожденная или плохо обусловленная $H$.** Решение линейной системы усиливает ошибки округления, и шаг может оказаться огромным.
\end{itemize}

\problemhead{B2}{Притяжение к максимуму}{Самостоятельно}
\textbf{Условие.} Примените полный метод Ньютона к $f(x)=\tfrac{1}{4}x^4-\tfrac{1}{2}x^2$ из $x_0=0{,}4$. Запишите итерацию как $x^+=\varphi(x)$ и найдите $x_1$, $x_2$. Куда сходится метод и что это за точка? Найдите наибольшее $a$, при котором любая начальная точка с $0<|x_0|< a$ притягивается к ней.

\textbf{Решение.}
\begingroup\sloppy
\begin{enumerate}[label=\textbf{\arabic*.}, leftmargin=*]
  \item $\varphi(x)=x-{x^3-x\over 3x^2-1}={2x^3\over 3x^2-1}$. Из $x_0=0{,}4$: $x_1=-0{,}2462$, $x_2=0{,}0365$.
  \item Итерации идут в $0$, где $f'(0)=0$ и $f''(0)=-1$: это локальный максимум. Метод Ньютона ищет $f'=0$ и не отличает максимум от минимума; возле $0$ $\varphi(x)\approx-2x^3$, так что сходимость даже кубическая.
  \item $\varphi(x)=-x$ даёт $2x^2=1-3x^2$, $x^2=1/5$. При $0<|x|< a$, где $a=1/\sqrt5=0{,}4472$, выполнено $|\varphi(x)|<|x|$, и итерации стягиваются к максимуму; при $|x|=a$ они зацикливаются.
\end{enumerate}
\endgroup
\begin{answerbox}
$x_1=-0{,}2462$, $x_2=0{,}0365$; сходимость к максимуму $x=0$ при $0<|x_0|<0{,}4472$.
\end{answerbox}

\begin{definitionbox}{Определение.}
\textbf{Демпфированный метод Ньютона} перебирает $\alpha_k=1,\beta,\beta^2,\dots$, пока не выполнится
\handoutdisplay{
f(x_k+\alpha_kp_k)\le f(x_k)+c\,\alpha_k\nabla f(x_k)^\top p_k,
}
 и полагает $x_{k+1}=x_k+\alpha_kp_k$. Ему нужно, чтобы $p_k$ было направлением спуска.
\end{definitionbox}

\begin{theorembox}{Утверждение.}
Если итерации демпфированного метода Ньютона сходятся к $x^*$ с $\nabla^2 f(x^*)\succ0$ и $c<\tfrac{1}{2}$, то для всех больших $k$ принимается $\alpha_k=1$, и сходимость квадратичная.
\end{theorembox}

\begin{proofbox}
Возле $x^*$ $f(x+p_N)-f(x)=-\tfrac{1}{2}g^\top H^{-1}g+o(\|g\|^2)$, а условие Армихо требует не больше $-c\,g^\top H^{-1}g$. Поскольку $c<\tfrac{1}{2}$, первое в итоге оказывается меньше второго.
\end{proofbox}

\problemhead{A1}{Ньютон на функции Розенброка: слишком длинный шаг}{У доски}
\textbf{Условие.} Для $f(x,y)=a(y-x^2)^2+(1-x)^2$ при $a=100$ начните из $x_0=(0;\ 0)$. Сделайте два полных шага Ньютона и найдите $f$ после каждого. Затем из $x_0$ запустите демпфированный метод: делите $\alpha$ пополам, начиная с $1$, пока не выполнится условие Армихо с $c=0{,}0001$. Какое $\alpha$ будет принято и чему там равна $f$?

\textbf{Решение.}
\begingroup\sloppy
\begin{enumerate}[label=\textbf{\arabic*.}, leftmargin=*]
  \item В $x_0$: $\nabla f=(-2;\ 0)$ и $\nabla^2 f=\begin{pmatrix} 2 & 0 \\ 0 & 200 \end{pmatrix}$. Решение $\nabla^2 f\,p=-\nabla f$ даёт $p_0=(1;\ 0)$, поэтому $x_1=(1;\ 0)$.
  \item Но $f(x_1)=100$, тогда как $f(x_0)=1$. Квадратичная модель в $x_0$ верна до второго порядка, однако её минимум лежит далеко за пределами области, где она описывает $f$: при $x=1$ дно долины $y=x^2$ находится на высоте один, а шаг остался на $y=0$.
  \item Из $x_1$: $\nabla f=(400;\ -200)$, $\nabla^2 f=\begin{pmatrix} 1202 & -400 \\ -400 & 200 \end{pmatrix}$, шаг Ньютона $p_1=(0;\ 1)$. Значит, $x_2=(1;\ 1)$ — минимум, $f(x_2)=0$. Полный метод туда приходит, но через точку, которая в сто раз хуже начальной.
  \item Демпфированный метод сохраняет направление $p_0$ и проверяет $f(x_0+\alpha p_0)\le f(x_0)+c\alpha\nabla f^\top p_0$, где $\nabla f^\top p_0=-2$. Для $\alpha=(1;\ 0{,}5;\ 0{,}25)$ значения равны $f=(100;\ 6{,}5;\ 0{,}9531)$: первые два не проходят, принимается $\alpha=0{,}25$ в точке $x=(0{,}25;\ 0)$.
\end{enumerate}
\endgroup
\begin{answerbox}
$f(x_1)=100$, $x_2=(1;\ 1)$; с демпфированием $\alpha=0{,}25$, пробные значения $f=(100;\ 6{,}5;\ 0{,}9531)$, последнее принято.
\end{answerbox}

% =====================================================================
\section{Сдвиг гессиана}\label{sec:regularization}

\begin{definitionbox}{Определение.}
\textbf{Регуляризованный шаг Ньютона} решает $(H+\lambda I)p=-g$, где $\lambda\ge0$ выбрано так, что $H+\lambda I\succ0$, то есть $\lambda>-\lambda_{\min}(H)$.
\end{definitionbox}

\begin{theorembox}{Утверждение.}
Если $H+\lambda I\succ0$ и $g\ne0$, то $g^\top p=-g^\top(H+\lambda I)^{-1}g<0$. При $\lambda\to0$ шаг стремится к шагу Ньютона, если $H\succ0$; при $\lambda\to\infty$ $p\approx-\lambda^{-1}g$ — короткий градиентный шаг.
\end{theorembox}

На практике $\lambda$ начинают с нуля или с малого значения и умножают на константу до тех пор, пока $H+\lambda I$ не допустит разложения Холецкого; после удачного полного шага $\lambda$ снова уменьшают. В собственном базисе шаг равен $p_i=-g_i/(\lambda_i+\lambda)$: по направлению отрицательной кривизны с $\lambda_i+\lambda>0$ теперь идут вниз.

\problemhead{B1}{Направление Ньютона ведёт вверх}{У доски}
\textbf{Условие.} Пусть $f(x,y)=\tfrac{1}{2}x^2-y^2+\tfrac{1}{4}y^4$ в точке $(0;\ 0{,}5)$. Найдите $g$, $H$, направление Ньютона $p_N$ и $g^\top p_N$. Является ли $p_N$ направлением спуска? Найдите наименьший сдвиг, при котором $H+\lambda I$ положительно определена, решите $(H+\lambda I)p=-g$ при $\lambda=2$ и проверьте полный шаг по Армихо с $c=0{,}0001$.

\textbf{Решение.}
\begingroup\sloppy
\begin{enumerate}[label=\textbf{\arabic*.}, leftmargin=*]
  \item $g=(0;\ -0{,}875)$, $H=\begin{pmatrix} 1 & 0 \\ 0 & -1{,}25 \end{pmatrix}$: одно собственное значение отрицательно, так что $H$ знаконеопределён. Метод Ньютона даёт $p_N=-H^{-1}g=(0;\ -0{,}7)$ и $g^\top p_N=0{,}6125>0$.
  \item У $\phi(\alpha)=f(x+\alpha p_N)$ производная $\phi'(0)=g^\top p_N>0$: любой малый шаг увеличивает $f$, а полный шаг переводит её из $-0{,}2344$ в $-0{,}0396$. Дробление уменьшало бы $\alpha$ бесконечно. Менять нужно само направление.
  \item $H+\lambda I\succ0$ ровно при $\lambda>1{,}25$. При $\lambda=2$: $p=(0;\ 1{,}1667)$ и $g^\top p=-1{,}0208<0$. Теперь шаг идёт вдоль отрицательной кривизны, прочь от седла.
  \item Армихо при $\alpha=1$: $f(x+p)=-0{,}8488\le f(x)+c\,g^\top p=-0{,}2345$, значит, полный регуляризованный шаг принимается. Минимумы находятся в $y=\pm1{,}4142$, где $H\succ0$ и сдвиг можно убрать.
\end{enumerate}
\endgroup
\begin{answerbox}
$g^\top p_N=0{,}6125$: это не направление спуска. $\lambda>1{,}25$; при $\lambda=2$ шаг равен $(0;\ 1{,}1667)$, и $f$ уменьшается до $-0{,}8488$.
\end{answerbox}

\handoutfigure{0.95}{fig_saddle}{Полный метод Ньютона сходится в седло; регуляризованный демпфированный метод уходит от него вдоль отрицательной кривизны.}

\begin{notebox}{Доверительные области}
Методы доверительной области выбирают сначала длину шага, а потом сдвиг: $p$ минимизирует модель при $\|p\|\le\Delta$, а решение имеет вид $(H+\lambda I)p=-g$ с $\lambda\ge0$, определяемым $\Delta$. Отрицательную кривизну они обрабатывают без отдельной проверки.
\end{notebox}

% =====================================================================
\section{Методы Гаусса–Ньютона и Левенберга–Марквардта}\label{sec:least-squares}

Для $f(x)=\tfrac{1}{2}\|r(x)\|^2$ с невязками $r:\mathbb R^n\to\mathbb R^N$ и якобианом $J$
\handoutdisplay{
\nabla f=J^\top r,\qquad \nabla^2 f=J^\top J+\sum_i r_i\nabla^2 r_i.
}


\begin{definitionbox}{Определение.}
\textbf{Метод Гаусса–Ньютона} отбрасывает второе слагаемое и решает $J^\top Jp=-J^\top r$; это шаг Ньютона для линеаризованной невязки $r+Jp$. \textbf{Метод Левенберга–Марквардта} решает $(J^\top J+\lambda I)p=-J^\top r$.
\end{definitionbox}

\begin{whybox}
$J^\top J\succeq0$ требует только первых производных, а отброшенное слагаемое мало, когда невязки возле решения малы. Сдвиг $\lambda$ справляется с неполным рангом $J$ и сдвигает шаг от Гаусса–Ньютона к короткому градиентному шагу.
\end{whybox}

\problemhead{B3}{Гаусс–Ньютон, Ньютон и Левенберг–Марквардт}{Самостоятельно}
\textbf{Условие.} Пусть $f(\theta)=\tfrac{1}{2}\|r(\theta)\|^2$, $r(\theta)=(\theta^2-2,\ \theta-1)$. Из $\theta_0=1$ сделайте один шаг Гаусса–Ньютона, один шаг Ньютона с точной второй производной и один шаг Левенберга–Марквардта с $\lambda=1$. Укажите $\theta_1$ для каждого и сравните с точкой минимума.

\textbf{Решение.}
\begingroup\sloppy
\begin{enumerate}[label=\textbf{\arabic*.}, leftmargin=*]
  \item В $\theta_0$: $r=(-1;\ 0)$, $J=(2;\ 1)$, $J^\top J=5$, $J^\top r=-2$. Гаусс–Ньютон: $\theta_1=\theta_0-J^\top r/J^\top J=1{,}4$.
  \item Точная вторая производная добавляет $r_1r_1''=-2$, так что $f''=3$, и метод Ньютона даёт $\theta_1=1{,}6667$. Левенберг–Марквардт делит на $J^\top J+\lambda$: $\theta_1=1{,}3333$.
  \item $f'(\theta)=2\theta^3-3\theta-1=(\theta+1)(2\theta^2-2\theta-1)$, точка минимума $\theta^*=(1+\sqrt3)/2=1{,}366$. Ближе всех оказывается Гаусс–Ньютон: в $\theta_0$ слагаемое с невязкой уменьшило точную кривизну, и шаг Ньютона вышел слишком длинным.
\end{enumerate}
\endgroup
\begin{answerbox}
Гаусс–Ньютон $1{,}4$, Ньютон $1{,}6667$, Левенберг–Марквардт $1{,}3333$; $\theta^*=1{,}366$.
\end{answerbox}

% =====================================================================
\section{Стоимость кривизны}\label{sec:cost}

\begin{center}\small
\begin{tabular}{@{}>{\raggedright\arraybackslash\hspace{0pt}}p{0.31\linewidth}>{\raggedright\arraybackslash\hspace{0pt}}p{0.39\linewidth}>{\raggedright\arraybackslash\hspace{0pt}}p{0.2\linewidth}@{}}
\toprule
Операция & Стоимость & Память \\
\midrule
Построить гессиан & $O(n^2)$ производных; $O(Nn^2)$ для GLM & $O(n^2)$ \\
Разложение Холецкого & $\approx n^3/3$ & $O(n^2)$ \\
Решение с разложением & $O(n^2)$ & $O(n^2)$ \\
Произведение гессиана на вектор & несколько градиентов & $O(n)$ дополнительно \\
Направление L-BFGS & $\approx 4mn$ & $2mn$ \\
\bottomrule
\end{tabular}
\end{center}

\begin{examplebox}
Логистическая функция потерь со штрафом $\ell_2$, $f(w)=N^{-1}\sum_i\bigl[\log(1+e^{z_i})-y_iz_i\bigr]+\tfrac{\gamma}{2}\|w\|^2$, $z=Xw$:
\handoutdisplay{
\nabla f=N^{-1}X^\top(\sigma(z)-y)+\gamma w,\quad Hv=N^{-1}X^\top\bigl(d\odot Xv\bigr)+\gamma v,
}
 где $d_i=\sigma(z_i)(1-\sigma(z_i))$. Произведение стоит двух проходов по $X$ и не строит $H$.
\end{examplebox}

\textbf{Newton-CG} приближённо решает $Hp=-g$ с помощью CG, которому нужны только произведения $Hv$. Ранняя остановка CG или остановка при отрицательной кривизне $v^\top Hv\le0$ всё равно даёт направление спуска.

\problemhead{C3}{Сколько стоит гессиан}{Самостоятельно}
\textbf{Условие.} Логистическая регрессия с $N=100000$ объектами и $n=20000$ признаками. Построение $H=N^{-1}X^\top DX$ стоит около $Nn^2$ умножений со сложением, одно произведение $Hv=N^{-1}X^\top(d\odot Xv)$ — около $2Nn$. Во сколько раз решение системы методом Newton-CG с $k=16$ произведениями дешевле построения $H$? Сколько гигабайт занимает плотная $H$ при $8$ байтах на элемент и сколько мегабайт занимают $m=5$ пар L-BFGS?

\textbf{Решение.}
\begingroup\sloppy
\begin{enumerate}[label=\textbf{\arabic*.}, leftmargin=*]
  \item $Nn^2/(2kNn)=n/(2k)=625$: число объектов сокращается, а выигрыш растёт вместе с числом признаков.
  \item Плотная $H$ содержит $n^2$ элементов: $3{,}2$ ГБ ещё до какого-либо разложения. L-BFGS хранит $2m$ векторов длины $n$: $1{,}6$ МБ.
  \item Разложение Холецкого добавило бы около $n^3/3$ операций. При $n$ порядка десятков тысяч метод Ньютона возможен только без матрицы, через $Hv$ и сопряжённые градиенты.
\end{enumerate}
\endgroup
\begin{answerbox}
Отношение $625$; плотная $H$ — $3{,}2$ ГБ; L-BFGS — $1{,}6$ МБ.
\end{answerbox}

\handoutfigure{0.95}{fig_cost}{Норма градиента по итерациям и по проходам по данным для градиентного спуска, демпфированного Ньютона, BFGS и L-BFGS.}

% =====================================================================
\section{Секущие пары и BFGS}\label{sec:bfgs}

По формуле Ньютона–Лейбница
\handoutdisplay{
y_k=\Bigl[\int_0^1\nabla^2 f(x_k+ts_k)\,dt\Bigr]s_k:
}
 $y_k$ — действие усреднённого гессиана на $s_k$. Квазиньютоновский метод хранит матрицу, которая его воспроизводит.

\begin{definitionbox}{Определение.}
\textbf{Секущие условия} — это $B_{k+1}s_k=y_k$ и $G_{k+1}y_k=s_k$. \textbf{Обновление BFGS} в форме гессиана и в обратной форме:
\handoutdisplay{
B_{k+1}=B_k-{B_ks_ks_k^\top B_k\over s_k^\top B_ks_k}+{y_ky_k^\top\over y_k^\top s_k},\qquad G_{k+1}=(I-\rho_ks_ky_k^\top)G_k(I-\rho_ky_ks_k^\top)+\rho_ks_ks_k^\top,
}
 где $\rho_k=1/(y_k^\top s_k)$. Направление $p_k=-G_kg_k$, с линейным поиском.
\end{definitionbox}

\begin{notebox}{Почему именно эта формула}
Одно векторное уравнение не определяет матрицу. BFGS — симметричное обновление, ближайшее к $G_k$ в некоторой взвешенной норме среди удовлетворяющих секущему условию; оно меняет $G$ на матрицу ранга не выше двух. DFP — та же конструкция, в которой $B$ и $G$ поменялись ролями.
\end{notebox}

\begin{theorembox}{Теорема.}
$G_{k+1}y_k=s_k$. Если $G_k\succ0$ и $y_k^\top s_k>0$, то $G_{k+1}\succ0$.
\end{theorembox}

\begin{proofbox}
$(I-\rho ys^\top)y=y-\rho y(s^\top y)=0$, поэтому первое слагаемое обращается в ноль на $y$, а $\rho ss^\top y=s$. Для $z\ne0$ положим $V=I-\rho sy^\top$: $z^\top G_{k+1}z=(V^\top z)^\top G_k(V^\top z)+\rho(s^\top z)^2\ge0$. Если бы это было нулём, то $s^\top z=0$, откуда $V^\top z=z-\rho y(s^\top z)=z$ и $z^\top G_kz=0$, то есть $z=0$.
\end{proofbox}

\problemhead{C1}{BFGS на квадратичной функции вручную}{У доски}
\textbf{Условие.} Пусть $f(x)=\tfrac{1}{2}x^\top Qx$, $Q=\begin{pmatrix} 2 & 1 \\ 1 & 1 \end{pmatrix}$, $x_0=(-2;\ 2)$, $G_0=I$, линейный поиск точный. Сделайте первый шаг, составьте $s$, $y$ и обратное обновление BFGS $G_1$. Проверьте секущее условие и положительную определённость. Сделайте второй шаг: чему равны $\alpha_1$ и $x_2$?

\textbf{Решение.}
\begingroup\sloppy
\begin{enumerate}[label=\textbf{\arabic*.}, leftmargin=*]
  \item $g_0=Qx_0=(-2;\ 0)$, $p_0=-G_0g_0=(2;\ 0)$, $\alpha_0=0{,}5$. Отсюда $s=(1;\ 0)$, $x_1=(-1;\ 2)$, $g_1=(0;\ 1)$ и $y=g_1-g_0=(2;\ 1)$; на квадратичной функции $y=Qs$ точно.
  \item $y^\top s=2>0$, $\rho=0{,}5$. При $V=I-\rho sy^\top$
\handoutdisplay{
G_1=VG_0V^\top+\rho ss^\top=\begin{pmatrix} 0{,}75 & -0{,}5 \\ -0{,}5 & 1 \end{pmatrix}.
}

  \item $G_1y=(1;\ 0)=s$. Матрица $G_1$ симметрична, её собственные значения $(0{,}3596;\ 1{,}3904)$ положительны. Обратная к ней $B_1=\begin{pmatrix} 2 & 1 \\ 1 & 1{,}5 \end{pmatrix}$ совпадает с $Q$ вдоль $s$, но пока не в других направлениях.
  \item $p_1=-G_1g_1=(0{,}5;\ -1)$, $g_1^\top p_1<0$, точный шаг $\alpha_1=2$: $x_2=(0;\ 0)=x^*$. При точном линейном поиске BFGS минимизирует $n$-мерную положительно определённую квадратичную функцию не более чем за $n$ шагов.
\end{enumerate}
\endgroup
\begin{answerbox}
$G_1=\begin{pmatrix} 0{,}75 & -0{,}5 \\ -0{,}5 & 1 \end{pmatrix}$, $G_1y=s$; $\alpha_1=2$, $x_2=(0;\ 0)$.
\end{answerbox}

\handoutfigure{0.85}{fig_bfgs}{Градиентный спуск, BFGS и метод Ньютона на функции Розенброка: зазор функции по итерациям.}

Возле невырожденного минимума BFGS с линейным поиском, который в итоге принимает $\alpha=1$, сходится сверхлинейно: $\|e_{k+1}\|/\|e_k\|\to0$. Квадратичной скорости Ньютона он не достигает, зато ему нужны только градиенты и $O(n^2)$ работы на шаг.

% =====================================================================
\section{Условие кривизны}\label{sec:curvature}

$y_k^\top s_k=s_k^\top\bigl[\int_0^1\nabla^2 f(x_k+ts_k)dt\bigr]s_k$ — средняя кривизна вдоль шага. Для $\mu$-сильно выпуклой $f$ она не меньше $\mu\|s_k\|^2>0$; для невыпуклой $f$ может иметь любой знак.

\begin{definitionbox}{Определение.}
\textbf{Условия Вульфа} для шага $\alpha$ вдоль направления спуска $p$ — это условие Армихо и условие кривизны $\nabla f(x+\alpha p)^\top p\ge c_2\nabla f(x)^\top p$, где $0< c< c_2<1$.
\end{definitionbox}

\begin{theorembox}{Утверждение.}
Если $\alpha$ удовлетворяет условию кривизны Вульфа вдоль направления спуска $p$, то $y^\top s>0$ для $s=\alpha p$.
\end{theorembox}

\begin{proofbox}
$y^\top s=\alpha\bigl(\nabla f(x+\alpha p)-\nabla f(x)\bigr)^\top p\ge\alpha(c_2-1)\nabla f(x)^\top p>0$, поскольку $c_2<1$ и $\nabla f(x)^\top p<0$.
\end{proofbox}

\problemhead{C2}{Условие кривизны}{Самостоятельно}
\textbf{Условие.} Для $f(x)=\tfrac{1}{4}x^4-\tfrac{1}{2}x^2$ возьмите точки $(0;\ 0{,}5;\ 1{,}5)$. Найдите $y^\top s$ для шага из первой точки во вторую; можно ли BFGS использовать эту пару? Для шага из второй точки в третью найдите секущую кривизну $B=y/s$ и квазиньютоновский шаг $x^+=x-f'(x)/B$ из третьей точки.

\textbf{Решение.}
\begingroup\sloppy
\begin{enumerate}[label=\textbf{\arabic*.}, leftmargin=*]
  \item Из первой точки во вторую: $s=0{,}5$, $y=-0{,}375$ и $y^\top s=-0{,}1875<0$. Обновление сделало бы $G$ отрицательной; BFGS должен пропустить эту пару.
  \item Из второй точки в третью: $s=1$, $y=2{,}25$, $B=y/s=2{,}25$. По теореме о среднем это среднее значение $f''$ на шаге, и оно положительно.
  \item Из третьей точки $x^+=x-f'(x)/B=0{,}6667$. Минимум находится в $1$; в одномерном случае BFGS — это метод секущих, который сходится сверхлинейно возле невырожденного минимума.
\end{enumerate}
\endgroup
\begin{answerbox}
$y^\top s=-0{,}1875$: пару нужно пропустить. $B=2{,}25$, $x^+=0{,}6667$.
\end{answerbox}

\begin{trapbox}
Одно дробление по Армихо не гарантирует $y^\top s>0$. Реализации пропускают обновление при $y^\top s\le\varepsilon\|y\|\|s\|$ или демпфируют $y$ в сторону $Bs$ (демпфирование Пауэлла). Без этой проверки одно обновление с $y^\top s<0$ может сделать $G$ знаконеопределённой, а следующее направление — направлением подъёма.
\end{trapbox}

% =====================================================================
\section{Ограниченная память и ограничения-отрезки}\label{sec:lbfgs}

Плотный BFGS хранит матрицу $n\times n$. \textbf{L-BFGS} хранит последние $m$ пар $(s_i,y_i)$, обычно $m$ от пяти до двадцати, и умножает $G_k$ на $g_k$, не строя её.

\begin{itemize}
  \item Положите $q=g_k$. Для пар от новой к старой: $a_i=\rho_is_i^\top q$, $q\leftarrow q-a_iy_i$.
  \item Масштабируйте: $r=\gamma_kq$, где $\gamma_k=s_{k-1}^\top y_{k-1}/y_{k-1}^\top y_{k-1}$ — оценка обратной кривизны.
  \item Для пар от старой к новой: $b_i=\rho_iy_i^\top r$, $r\leftarrow r+s_i(a_i-b_i)$.
  \item Верните $p_k=-r$: около $4mn$ операций и $2mn$ хранимых чисел.
\end{itemize}

Результат равен $-Gg$ для матрицы, полученной из $\gamma_kI$ последовательными обновлениями BFGS по хранимым $m$ парам. Старая кривизна забывается, и на задачах, где она была бы полезна, это стоит итераций.

\begin{notebox}{L-BFGS-B}
Для ограничений $l\le x\le u$ \textbf{L-BFGS-B} движется по пути проекции градиента до обобщённой точки Коши, фиксирует переменные на их границах, минимизирует модель с ограниченной памятью по свободным переменным и делает линейный поиск, не выходящий из бруса. Стационарность измеряется проекцией градиента.
\end{notebox}

\begin{trapbox}
Обрезка решения задачи без ограничений по брусу в общем случае не даёт условный минимум: если две переменные связаны, фиксация одной на границе сдвигает оптимум другой.
\end{trapbox}

\begin{center}\small
\begin{tabular}{@{}lll>{\raggedright\arraybackslash\hspace{0pt}}p{0.46\linewidth}@{}}
\toprule
Метод & Нужно & На шаг & Возле решения \\
\midrule
Градиентный спуск & $g$ & $O(n)$ & линейная, скорость задаёт $\kappa$ \\
Ньютон & $g$, $H$ & $O(n^3)$ & квадратичная \\
Newton-CG & $g$, $Hv$ & итерации CG & сверхлинейная при уменьшающемся допуске \\
BFGS & $g$ & $O(n^2)$ & сверхлинейная \\
L-BFGS & $g$ & $O(mn)$ & в общем случае линейная \\
\bottomrule
\end{tabular}
\end{center}

% =====================================================================
\section{Проверки реализации и упражнения}\label{sec:practice}

\begin{itemize}
  \item Решайте систему, а не обращайте матрицу; определяйте $H\not\succ0$ по неудавшемуся разложению Холецкого и сдвигайте.
  \item Проверяйте $g^\top p<0$ до линейного поиска и $y^\top s>0$ до обновления BFGS.
  \item Проверяйте производные: центральные разности для $g$, разность градиентов для $Hv$, для нескольких величин шага.
  \item Записывайте вызовы функции, градиента и гессиана рядом с итерациями и сравнивайте методы при одном правиле остановки.
  \item В конце проверяйте гессиан или наименьшую найденную кривизну: малый градиент может означать седло.
\end{itemize}

\begin{selfcheckbox}
Почему дробление шага не может исправить направление Ньютона с $g^\top p>0$?
\selfcheckanswer Потому что $f(x+\alpha p)=f(x)+\alpha g^\top p+o(\alpha)$ растёт при любом малом $\alpha$; условие Армихо не выполняется ни для одного из них.
\end{selfcheckbox}

\begin{selfcheckbox}
Какую информацию несёт одна секущая пара $(s,y)$?
\selfcheckanswer Действие среднего гессиана на отрезке на один вектор $s$ и ничего о других направлениях.
\end{selfcheckbox}

\begin{selfcheckbox}
Когда произведение гессиана на вектор полезнее самого гессиана?
\selfcheckanswer Когда $n$ велико: $Hv$ стоит несколько градиентов и $O(n)$ памяти, а методу сопряжённых градиентов нужны только такие произведения.
\end{selfcheckbox}

\end{document}
